\begin{recipe}{Stoofpeertjes}
\kicker[Dessert,Zoet,Nederlands]{Dessert · zoet · Nederlands}
\index[register]{Stoofpeertjes}
\index[register]{Peer}
\index[register]{Rode wijn}
\index[register]{Kaneel}

\meta{20 minuten + 2 tot 3 uur stoven \dvd Voor 6 personen \dvd Eén pan}
\marginimage{images/stoofpeertjes-hero}

\lettrine{S}{toofpeertjes} zijn een klassiek Hollands toetje: stevige stoofperen
die uren zachtjes garen in rode wijn met kaneel en kruidnagel, tot ze een
diepe rode kleur krijgen en vanzelf zacht worden. Het meeste werk laat je aan
de pan over. Lekker warm of koud als toetje, bijvoorbeeld met een schepje
vla, maar bij oma gingen ze net zo vaak gewoon als bijgerecht mee.

\blockrule{Ingrediënten}
\begin{ingredients}
  \ing{1 kg}{stoofperen (bij voorkeur Gieser Wildeman)}\index[register]{Peer}
  \ing{75--100 g}{kristalsuiker (of bruine basterdsuiker, voor een extra warme smaak)}
  \ing{½--1 tl}{kaneelpoeder}\index[register]{Kaneel}
  \ing{1 mespunt tot ¼ tl}{kruidnagelpoeder}
  \ing{250 ml}{rode wijn (donkerrood en fruitig, bijvoorbeeld merlot, pinot noir of shiraz)}\index[register]{Rode wijn}
  \ing{150--250 ml}{water, om aan te vullen}
\end{ingredients}

\blockrule{Bereiding}
\tip{Gebruik een pan die net groot genoeg is voor het aantal peren: hoe
minder ruimte eromheen, hoe minder vocht je nodig hebt om ze onder te
zetten.}
\begin{steps}
  \item Was de peren en schil ze rondom met een dunschiller. Laat het
        steeltje eraan zitten voor de klassieke look. Laat de peren heel, of
        halveer ze of snijd ze in kwarten en verwijder het klokhuis.
  \item Zet de peren rechtop in een hapjespan met deksel en strooi de suiker,
        het kaneelpoeder en het kruidnagelpoeder erover.
  \item Giet de rode wijn erbij en vul aan met water tot de peren net
        onderstaan. Roer kort en voorzichtig door, zodat het kaneel- en
        kruidnagelpoeder oplossen en niet aan het oppervlak klonteren.
  \item Breng het geheel aan de kook op middelhoog vuur. Zet zodra het kookt
        het vuur direct op de allerlaagste stand en doe de deksel op de pan.
  \item Laat de peren 2 tot 3 uur zachtjes stoven, tot ze zacht zijn en een
        diepe rode kleur hebben aangenomen. Hoe langer ze stoven, hoe zachter
        de peren en hoe dieper de kleur.
\end{steps}
\end{recipe}
